\subsection{$L_1$ Sampling Protocol with Secret Shared Output}  \label{apx:ssL1}
 
We give a more formal description of the functionality $\F_{L_1}^{ss}$. 

\bbox
{
\begin{center}
 $\F_{L_1}$:  Ideal functionality for two-party $L_1$ sampling
\end{center}

The functionality has the following parameter:
\BI
\item $n \in \NN$. The dimension of the input weight vectors $\vec{w}_1$ and $\vec{w}_2$.
\EI

The functionality proceeds as follows:
\BEN
\item Receive inputs $\vec w_1$ and $\vec w_2$ from $P_1$ and $P_2$
respectively. 
\item Sample $i \in [n]$ with probability $\frac{w_{1,i} +
w_{2,i}}{ \|\vec{w}_1 + \vec{w}_2\|_1 }$
\item Choose a random number $\pi$ consisting of $\lceil \log n \rceil$ bits.
\item Send $\pi$ to $P_1$ and $i \xor \pi$ to $P_2$. 
\EEN 

}

\medskip

We describe a protocol securely realizing
$\F_{L_1}^{ss}$ in the $(\F_\oslone, \F_\bcoin)$-hybrid. 

\floatname{algorithm}{Protocol}
\begin{algorithm}[h]
\textbf{Inputs:} Party $P_b$ has input $\vec{w}_b$. 

\BEN
\item Execute $\F_\oslone$ with $P_1$ as a sender with input $\vec{w}_1$
and $P_2$ as a receiver. Let $\langle i_1 \rangle$ be the secret share of the output
index. 

\item Execute $\F_\oslone$ with $P_2$ as a sender with input $\vec{w}_2$
and $P_1$ as a receiver. Let $\langle i_2 \rangle$ be the secret share of the output
index.

\item Execute $\F_\bcoin$ where $P_1$ has input $\|\vec{w}_1\|_1$ and
$P_2$ has input $\|\vec{w}_2\|_1$. Let $\langle b \rangle$ be the secret share of the
output bit. In addition, $P_1$ chooses random bits $\pi$.

\item Execute $\F_{2PC}$ for the following circuit: 
      \BEN
      \item Input: $\langle i_1 \rangle, \langle i_2 \rangle, \langle b \rangle, \pi$. 
      \item Compute $i = i_1 \cdot (1-b) + i_2 \cdot b$.
      \item Output $\pi$ to $P_1$ and $\pi \xor i$ to $P_2$.
      \EEN
\EEN

\caption{Protocol securely realizing $\F_{L_1}^{ss}$ in the $(\F_\oslone, \F_\bcoin)$-hybrid.}
\end{algorithm}

Security of the protocol can be shown similarly to Theorem~\ref{th:l1}.  
